\documentclass[11pt,a4paper]{article}
\newcommand{\intestazione}{Dimero di spin --- Schema dei quattro passi}
% =====================================================================
%  Preambolo comune ai documenti sul dimero (serie "che cosa abbiamo fatto")
%  Incluso con % =====================================================================
%  Preambolo comune ai documenti sul dimero (serie "che cosa abbiamo fatto")
%  Incluso con % =====================================================================
%  Preambolo comune ai documenti sul dimero (serie "che cosa abbiamo fatto")
%  Incluso con \input{preambolo_dimero} subito dopo \documentclass.
% =====================================================================
\usepackage[T1]{fontenc}
\usepackage[utf8]{inputenc}
\usepackage[main=italian,provide=*]{babel}
\usepackage{amsmath,amssymb,mathtools}
\usepackage{graphicx}
\usepackage{booktabs}
\usepackage{colortbl}
\usepackage{xcolor}
\usepackage{geometry}
\usepackage{placeins}
\usepackage{caption}
\usepackage{enumitem}
\usepackage{fancyhdr}
\usepackage{needspace}
\usepackage{tikz}
\usetikzlibrary{arrows.meta,positioning,shapes.geometric,calc,fit,backgrounds}
\usepackage{hyperref}
\hypersetup{colorlinks=true,linkcolor=blu,citecolor=blu,urlcolor=blu}

\geometry{margin=2.5cm}

% --- colori coerenti con quelli delle figure (palette Okabe-Ito)
\definecolor{blu}{HTML}{0072B2}
\definecolor{arancio}{HTML}{E69F00}
\definecolor{verdes}{HTML}{009E73}
\definecolor{rossos}{HTML}{D55E00}
\definecolor{violas}{HTML}{CC79A7}
\definecolor{grigetto}{HTML}{E8E8E8}

\captionsetup{font=small,labelfont=bf,width=0.92\textwidth}

\pagestyle{fancy}
\fancyhf{}
\fancyhead[L]{\small\itshape\intestazione}
\fancyhead[R]{\small\thepage}
\renewcommand{\headrulewidth}{0.3pt}

% --- notazione
\newcommand{\ket}[1]{\lvert #1\rangle}
\newcommand{\bra}[1]{\langle #1 \rvert}
\newcommand{\media}[1]{\langle #1 \rangle}
\newcommand{\vs}[1]{\mathbf{s}_{#1}}

% --- riquadro "in una riga": il messaggio da portare a casa
\newcommand{\messaggio}[1]{%
  \begin{center}
  \begin{minipage}{0.93\textwidth}
  \setlength{\fboxsep}{9pt}%
  \colorbox{grigetto}{\begin{minipage}{\dimexpr\textwidth-2\fboxsep}
  \small\textbf{In una riga.}\enspace #1
  \end{minipage}}
  \end{minipage}
  \end{center}
}

% --- riquadro "come lo sappiamo": la verifica numerica dietro un'affermazione
\newcommand{\verifica}[1]{%
  \begin{center}
  \begin{minipage}{0.93\textwidth}
  \small\textcolor{verdes}{\rule{2.2em}{1.1pt}}\enspace
  \textcolor{verdes}{\textbf{Come lo sappiamo.}}\enspace #1
  \end{minipage}
  \end{center}
  \medskip
}
 subito dopo \documentclass.
% =====================================================================
\usepackage[T1]{fontenc}
\usepackage[utf8]{inputenc}
\usepackage[main=italian,provide=*]{babel}
\usepackage{amsmath,amssymb,mathtools}
\usepackage{graphicx}
\usepackage{booktabs}
\usepackage{colortbl}
\usepackage{xcolor}
\usepackage{geometry}
\usepackage{placeins}
\usepackage{caption}
\usepackage{enumitem}
\usepackage{fancyhdr}
\usepackage{needspace}
\usepackage{tikz}
\usetikzlibrary{arrows.meta,positioning,shapes.geometric,calc,fit,backgrounds}
\usepackage{hyperref}
\hypersetup{colorlinks=true,linkcolor=blu,citecolor=blu,urlcolor=blu}

\geometry{margin=2.5cm}

% --- colori coerenti con quelli delle figure (palette Okabe-Ito)
\definecolor{blu}{HTML}{0072B2}
\definecolor{arancio}{HTML}{E69F00}
\definecolor{verdes}{HTML}{009E73}
\definecolor{rossos}{HTML}{D55E00}
\definecolor{violas}{HTML}{CC79A7}
\definecolor{grigetto}{HTML}{E8E8E8}

\captionsetup{font=small,labelfont=bf,width=0.92\textwidth}

\pagestyle{fancy}
\fancyhf{}
\fancyhead[L]{\small\itshape\intestazione}
\fancyhead[R]{\small\thepage}
\renewcommand{\headrulewidth}{0.3pt}

% --- notazione
\newcommand{\ket}[1]{\lvert #1\rangle}
\newcommand{\bra}[1]{\langle #1 \rvert}
\newcommand{\media}[1]{\langle #1 \rangle}
\newcommand{\vs}[1]{\mathbf{s}_{#1}}

% --- riquadro "in una riga": il messaggio da portare a casa
\newcommand{\messaggio}[1]{%
  \begin{center}
  \begin{minipage}{0.93\textwidth}
  \setlength{\fboxsep}{9pt}%
  \colorbox{grigetto}{\begin{minipage}{\dimexpr\textwidth-2\fboxsep}
  \small\textbf{In una riga.}\enspace #1
  \end{minipage}}
  \end{minipage}
  \end{center}
}

% --- riquadro "come lo sappiamo": la verifica numerica dietro un'affermazione
\newcommand{\verifica}[1]{%
  \begin{center}
  \begin{minipage}{0.93\textwidth}
  \small\textcolor{verdes}{\rule{2.2em}{1.1pt}}\enspace
  \textcolor{verdes}{\textbf{Come lo sappiamo.}}\enspace #1
  \end{minipage}
  \end{center}
  \medskip
}
 subito dopo \documentclass.
% =====================================================================
\usepackage[T1]{fontenc}
\usepackage[utf8]{inputenc}
\usepackage[main=italian,provide=*]{babel}
\usepackage{amsmath,amssymb,mathtools}
\usepackage{graphicx}
\usepackage{booktabs}
\usepackage{colortbl}
\usepackage{xcolor}
\usepackage{geometry}
\usepackage{placeins}
\usepackage{caption}
\usepackage{enumitem}
\usepackage{fancyhdr}
\usepackage{needspace}
\usepackage{tikz}
\usetikzlibrary{arrows.meta,positioning,shapes.geometric,calc,fit,backgrounds}
\usepackage{hyperref}
\hypersetup{colorlinks=true,linkcolor=blu,citecolor=blu,urlcolor=blu}

\geometry{margin=2.5cm}

% --- colori coerenti con quelli delle figure (palette Okabe-Ito)
\definecolor{blu}{HTML}{0072B2}
\definecolor{arancio}{HTML}{E69F00}
\definecolor{verdes}{HTML}{009E73}
\definecolor{rossos}{HTML}{D55E00}
\definecolor{violas}{HTML}{CC79A7}
\definecolor{grigetto}{HTML}{E8E8E8}

\captionsetup{font=small,labelfont=bf,width=0.92\textwidth}

\pagestyle{fancy}
\fancyhf{}
\fancyhead[L]{\small\itshape\intestazione}
\fancyhead[R]{\small\thepage}
\renewcommand{\headrulewidth}{0.3pt}

% --- notazione
\newcommand{\ket}[1]{\lvert #1\rangle}
\newcommand{\bra}[1]{\langle #1 \rvert}
\newcommand{\media}[1]{\langle #1 \rangle}
\newcommand{\vs}[1]{\mathbf{s}_{#1}}

% --- riquadro "in una riga": il messaggio da portare a casa
\newcommand{\messaggio}[1]{%
  \begin{center}
  \begin{minipage}{0.93\textwidth}
  \setlength{\fboxsep}{9pt}%
  \colorbox{grigetto}{\begin{minipage}{\dimexpr\textwidth-2\fboxsep}
  \small\textbf{In una riga.}\enspace #1
  \end{minipage}}
  \end{minipage}
  \end{center}
}

% --- riquadro "come lo sappiamo": la verifica numerica dietro un'affermazione
\newcommand{\verifica}[1]{%
  \begin{center}
  \begin{minipage}{0.93\textwidth}
  \small\textcolor{verdes}{\rule{2.2em}{1.1pt}}\enspace
  \textcolor{verdes}{\textbf{Come lo sappiamo.}}\enspace #1
  \end{minipage}
  \end{center}
  \medskip
}

\usepackage{pdflscape}
\usepackage{array}

\title{\bfseries Schema riassuntivo\\[2pt]
\large Punti di lavoro, stati iniziali e collegamenti nei quattro passi}
\author{Samuele Schiavi \\ \small Tesi triennale in Fisica, Università di Parma
--- Relatore: Prof.\ Alessandro Chiesa}
\date{}

\begin{document}
\maketitle
\thispagestyle{fancy}

Questo documento non introduce fisica nuova: raccoglie in un unico posto,
per ciascuno dei quattro passi già svolti, il punto di lavoro usato, lo
stato iniziale, il risultato chiave e il collegamento con quanto viene
dopo --- pensato per essere consultato, non letto in ordine.

\section*{Notazione: due basi diverse, stesso spazio}

Da qui in avanti compaiono due notazioni che si assomigliano ma
\textbf{non sono la stessa cosa} --- vanno tenute distinte fin dall'inizio,
non chiarite a parole ogni volta che ricorrono.

\begin{center}
\renewcommand{\arraystretch}{1.5}
\begin{tabular}{@{}cccc@{}}
\toprule
\rowcolor{grigetto}
\textbf{stato $\ket{S,M}$} & \textbf{nome} & \textbf{forma esplicita} & \textbf{stato di registro} \\
\midrule
$\ket{1,1}$ & tripletto, $M{=}{+}1$ & $\uparrow\uparrow$ & $\ket{00}$ \\
$\ket{1,0}$ & tripletto, $M{=}0$ & $(\uparrow\downarrow+\downarrow\uparrow)/\sqrt2$ & $(\ket{01}+\ket{10})/\sqrt2$ \\
$\ket{0,0}$ & singoletto & $(\uparrow\downarrow-\downarrow\uparrow)/\sqrt2$ & $(\ket{01}-\ket{10})/\sqrt2$ \\
$\ket{1,-1}$ & tripletto, $M{=}{-}1$ & $\downarrow\downarrow$ & $\ket{11}$ \\
\bottomrule
\end{tabular}
\end{center}

\textbf{Colonna di sinistra} ($\ket{S,M}$, sempre con la virgola): la base
di \emph{spin totale}, usata quando il significato fisico conta --- il
singoletto, il tripletto e le sue tre componenti.
\textbf{Colonna di destra} ($\ket{00}$ e simili, sempre senza virgola): la
base di \emph{registro} --- lo stato letterale dei due qubit, quello che
un circuito manipola davvero, prima di qualunque significato fisico
attribuito a posteriori.

\messaggio{La virgola non è un dettaglio tipografico: $\ket{1,1}$ (con la
virgola) è un'etichetta di spin totale, un'affermazione fisica.
$\ket{00}$ (senza virgola) è solo il nome di due bit messi a zero --- un
fatto sul registro, non ancora una affermazione fisica su cosa
rappresenti. Da qui in poi, ogni stato di registro citato nel testo è
accompagnato dalla parola ``registro'', non lasciato alla sola notazione.}

\section{Tabella d'insieme}

\renewcommand{\arraystretch}{1.35}
\begin{landscape}
\thispagestyle{empty}
\begin{center}
\small
\setlength{\tabcolsep}{4pt}
\begin{tabular}{@{}>{\bfseries\raggedright}p{2.7cm} p{5.15cm} p{5.15cm} p{5.15cm} p{5.15cm}@{}}
\toprule
\rowcolor{grigetto}
 & \textbf{Il sistema} & \textbf{Il VQE} & \textbf{La dinamica (Trotter)} & \textbf{I correlatori} \\
\midrule
Obiettivo &
Spettro esatto del dimero; ruolo del termine DM sull'incrocio dei livelli &
Trovare il fondamentale con un circuito; capire quale ansatz funziona, e
perché &
Evoluzione temporale reale $e^{-iHt}$; quantificare l'errore di Trotter &
Correlatori dinamici $C_{ij}^{\alpha\beta}(t)$ con un circuito ad ancilla \\
\addlinespace
Punto/i di lavoro &
sweep in $b/J$; confronto $D=0$ vs $D/J=0.2$ &
sweep in $b/J$, $D=0$ e $D/J=0.2$; più il punto singolo ``test 2''
($b/J{=}0.35$, $D/J{=}0.80$) &
$b/J=-0.18$, $D/J=1$ (DM forte, $5\times$ quello del VQE principale) &
``test 2'' ($b/J{=}0.35$, $D/J{=}0.80$) --- stesso punto singolo del VQE \\
\addlinespace
Stato iniziale &
nessuno (solo diagonalizzazione) &
$\ket{00}$ per costruzione (registro di default). L'ansatz generico è
tutto parametrizzato, senza porte fisse. Quello PMA parte invece da un
punto scelto apposta: singoletto oppure registro $\ket{11}$ per l'ansatz
minimo (dipende dal settore, $b/J$ sopra o sotto l'incrocio); sempre
registro $\ket{10}$ per quello esteso, senza scelta --- poi vengono le
porte parametrizzate &
il registro $\ket{00}$ stesso per il grosso del passo (coincide con
$\ket{1,1}$, nessuna porta applicata); nella sezione finale lo stesso
registro $\ket{00}$ è invece \emph{trasformato dalle porte} PMA
base/esteso (\emph{ri-ottimizzate a questo stesso punto}) in
un'approssimazione del fondamentale, usata al posto di $\ket{1,1}$ &
fondamentale esatto per il grosso del passo; PMA base/esteso
ri-ottimizzati allo stesso punto, solo nella sezione finale \\
\addlinespace
Risultato chiave &
incrocio esatto a $b/J=2$ con $D=0$ (gap$=0$); con $D/J=0.2$ si apre
(gap minimo $0.565$) &
ansatz base (CNOT--$R_y$--CNOT) crolla vicino all'incrocio se $D\neq0$
($\mathcal F\to0.70$); risolto con RBS $+\,2$ rotazioni ($\mathcal F=1$
sempre) &
errore $\sim1/N$ (operatori), $\sim1/N^2$ (fidelity); con lo stato VQE
minimo l'errore di preparazione ($0.327$) domina $10\times$ quello di
Trotter ($0.024$) &
stessa gerarchia di errori vista nella dinamica, ora sul correlatore; lo
scan sistematico (circuito vero) mostra un fattore $\sim12\times$ fra
combinazione di spin più protetta e più fragile con il PMA base; con
l'esteso, errore di Trotter reale ma piccolo ($0.011$--$0.032$) \\
\addlinespace
Collegamento &
motiva il passo successivo: l'ansatz dovrà fare i conti con una
simmetria che, con $D\neq0$, non c'è più &
fornisce la \emph{struttura di circuito} (PMA base/esteso) --- non uno
stato --- riusata indipendentemente nei due passi seguenti, ciascuno al
proprio punto &
introduce il tema ``due errori distinti, non intercambiabili'', ripreso
tale e quale nel passo successivo &
chiude la serie: stessa disciplina (Esatto/Vero dichiarati esplicitamente,
metodi confrontati non sostituiti, ansatz VQE verificato) applicata a una
quarta grandezza fisica \\
\bottomrule
\end{tabular}
\end{center}
\end{landscape}

\clearpage

\section{Il sistema}

Punto di partenza puramente teorico: diagonalizzazione esatta della
matrice $4\times4$ del dimero, nessun circuito. Il risultato che conta per
tutto il resto del lavoro non è lo spettro in sé, ma la scoperta che il
termine DM \textbf{rompe una simmetria esplicita}: senza DM,
l'Hamiltoniana conserva la magnetizzazione totale $M$, e i livelli con $M$
diverso possono incrociarsi senza respingersi (incrocio protetto, non una
singolarità --- il gap resta perfettamente definito, si annulla senza che
nulla diverga). Con il DM acceso, quella protezione sparisce: l'incrocio
si apre in un gap minimo finito.

\messaggio{Non è un dettaglio isolato: è la ragione fisica per cui, nel
passo successivo, un ansatz costruito per rispettare la conservazione di
$M$ funziona perfettamente senza DM e fallisce con il DM acceso ---
un'unica idea che attraversa tutto il lavoro.}

\section{Il VQE}

Circuito parametrico ottimizzato per trovare il fondamentale, confrontato
con la diagonalizzazione esatta. Il confronto centrale è fra due ansatz:
uno che rispetta la conservazione di $M$ (pochi parametri, blocco
CNOT--$R_y$--CNOT o, nella versione raffinata, RBS) e uno generico che non
rispetta nessuna simmetria particolare (sei parametri, sempre corretto ma
più costoso). Senza DM il primo vince nettamente: stessa fidelity del
generico, un sesto dei parametri. Con il DM acceso, vicino all'incrocio,
il primo crolla ($\mathcal F\approx0.70$) --- non per pochi parametri, ma
perché è confinato in un settore di $M$ che il vero fondamentale, con il
DM acceso, non abita più per intero. Aggiungendo due rotazioni che rompono
volutamente quella conservazione (ansatz esteso, RBS $+\,2$ parametri), la
fidelity torna a $1$ ovunque.

Oltre alla scansione principale, un secondo controllo puntuale: al punto
``test 2'' ($b/J=0.35$, $D/J=0.80$), scelto apposta perché il
fondamentale sia una sovrapposizione genuina e il DM sia ben acceso, la
stessa storia si conferma --- $\mathcal F=0.9402$ per l'ansatz base,
$\mathcal F=1$ (a dodici cifre) per quello esteso. Non è un punto
qualsiasi: è quello che il passo dei correlatori userà più avanti come
suo punto di lavoro principale, proprio perché già verificato qui.

Un fatto che vale la pena rendere esplicito, richiamando la tabella
d'apertura: le porte fisse dell'ansatz base, per $b/J$ sotto
l'incrocio, trasformano il registro esattamente nel singoletto
$\ket{0,0}$ --- e questo \emph{non} è un dettaglio qualsiasi: è
\textbf{esattamente} il fondamentale vero a $D=0$, verificato per
confronto diretto con la diagonalizzazione. La porta parametrizzata che
segue non deve quindi costruire il fondamentale da zero: deve solo
correggere la deviazione introdotta dal DM, allontanandosi da quel punto
di partenza già esatto quando $D=0$. È proprio per questo che un solo
parametro basta quando il DM è spento, e non basta più quando è acceso:
la correzione da fare non è più una semplice rotazione dentro lo stesso
settore.

L'ansatz esteso non funziona allo stesso modo, ed è un punto facile da
fraintendere: la sua porta fissa (una singola $X$, sempre la stessa,
prima di RBS) prepara $\ket{10}$ --- uno stato \emph{prodotto}, non
entangled, non il singoletto (sovrapposizione col vero singoletto:
$1/\sqrt2$, non $1$). Non serve che lo sia: RBS esplora \emph{da sola}
l'intero sottospazio $\{\ket{01},\ket{10}\}$, e il singoletto stesso è
uno dei punti che raggiunge, esattamente a $\varphi=\pi/4$ --- verificato
che $\text{RBS}(\pi/4)\ket{10}$ coincide col singoletto a precisione
macchina. L'entanglement qui non è un dato di partenza: è il
\emph{risultato} della stessa porta che l'ottimizzatore addestra.

\begin{center}
\small
\renewcommand{\arraystretch}{1.3}
\begin{tabular}{@{}p{2.6cm}p{4.9cm}p{4.9cm}@{}}
\toprule
\rowcolor{grigetto}
 & \textbf{ansatz base} & \textbf{ansatz esteso} \\
\midrule
Porte fisse &
$X,H,\text{CNOT},X$ (o $X,X$), due rami &
$X$ soltanto, sempre la stessa \\
Stato preparato &
singoletto o $\ket{11}$ --- \textbf{già entangled} &
$\ket{10}$ --- prodotto, \textbf{non entangled} \\
Perché &
il blocco dopo non sa cambiare settore: deve partire già in quello giusto &
RBS sa muoversi da sola: il punto di partenza preciso non conta \\
\bottomrule
\end{tabular}
\end{center}

\messaggio{Il criterio di progetto che ne esce: un ansatz va costruito
sulle simmetrie che l'Hamiltoniana ha davvero, non su quelle che sarebbe
comodo avesse. Non è un risultato che resta qui --- è la struttura di
circuito riusata, ri-ottimizzata da capo, nei due passi seguenti.}

\section{La dinamica (Trotter)}

Evoluzione temporale reale, $e^{-iHt}$ scomposto in passi di Trotter. Punto
di lavoro diverso da quello del VQE principale ($D/J=1$, cinque volte più
forte), scelto non per continuità con il VQE ma per una ragione propria:
serve una dinamica con più frequenze visibili insieme, altrimenti l'errore
di Trotter è quasi invisibile.

Il passo si scompone in \textbf{tre confronti distinti}, non due, tutti
allo stesso punto di lavoro ($b/J=-0.18$, $D/J=1$):
\begin{enumerate}[label=(\arabic*),leftmargin=2.2em,itemsep=4pt]
\item $\ket{1,1}\to$ Trotter, da solo --- quantifica l'errore di Trotter al
  crescere di $N$; autosufficiente, non serve nient'altro perché conti già
  come verifica del metodo (vedi sotto);
\item PMA base (parte sempre dal registro $\ket{00}$, ri-ottimizzato qui)
  $\to$ Trotter --- mostra l'effetto di una preparazione imperfetta
  ($\mathcal F=0.9402$);
\item PMA esteso (parte sempre dal registro $\ket{00}$, ri-ottimizzato qui)
  $\to$ Trotter, come controllo --- preparazione quasi perfetta
  ($\mathcal F\approx1$), resta solo l'errore di Trotter.
\end{enumerate}
Il primo è il grosso del passo; il secondo e il terzo compaiono insieme,
affiancati in due pannelli, nella sezione finale.

\subsection*{Il primo confronto: perché $\ket{1,1}$}

Lo stato di partenza usato per il primo confronto è $\ket{\uparrow\uparrow}$
--- equivalentemente $\ket{1,1}$, la componente $S{=}1,M{=}1$ del
tripletto, che per due spin-$1/2$ coincide esattamente con il prodotto
$\ket{\uparrow\uparrow}$, senza bisogno di sovrapposizioni --- non il
fondamentale.

La scelta non è un dettaglio implementativo: è essa stessa la condizione
per cui questo confronto conti come una \textbf{verifica} del metodo, e
non solo come una sua dimostrazione. Uno stato che fosse già autostato dei
singoli blocchi $H_1$, $H_2$ evolverebbe in modo banale sotto ciascuno di
essi, e non mostrerebbe l'errore di Trotter anche quando presente --- un
test di convergenza fatto su uno stato simile darebbe un risultato ottimo,
ma privo di significato, perché non starebbe davvero mettendo alla prova
l'approssimazione. $\ket{1,1}$ non è autostato di nessuno dei due blocchi a
questo punto di lavoro: la dinamica che ne segue è genuinamente
multi-frequenza, ed è proprio questo a rendere il confronto al crescere di
$N$ una verifica reale, non un caso scelto per apparire bene.

\subsection*{Il secondo e terzo confronto: PMA base ed esteso}

Qui si introduce lo stato del VQE --- non quello ottimizzato nel passo
precedente a $D/J=0.2$, ma la \emph{stessa struttura di circuito} (PMA
base, PMA esteso), che parte sempre dal registro $\ket{00}$ come ogni ansatz, e
viene \textbf{ri-ottimizzata da zero} a questo nuovo punto ($b/J=-0.18$,
$D/J=1$). Il risultato: con l'ansatz base, l'errore di preparazione
($0.327$) è più di dieci volte quello di Trotter ($0.024$) --- aumentare
$N$ non aiuta finché la preparazione resta quella. Con l'ansatz esteso
l'errore di preparazione crolla a zero per costruzione, e quel che resta
($0.033$) è solo Trotter.

Restare sullo stesso punto, invece di tornare a quello originale del VQE
($D/J=0.2$), non è un dettaglio secondario: è la condizione che rende il
confronto interpretabile. Lo scopo di questi due confronti è isolare un
solo effetto --- cosa cambia se la preparazione è imperfetta, a parità di
tutto il resto --- e per farlo va cambiata una sola variabile alla volta.
Cambiare anche il punto di lavoro mescolerebbe due effetti diversi
(preparazione imperfetta \emph{e} fisica diversa: frequenze, gap e
ricchezza della dinamica cambierebbero insieme), rendendo impossibile
dire a quale dei due attribuire uno scarto osservato. Per questo il VQE
rientra qui come \emph{metodo} di preparazione --- la stessa struttura di
circuito, ri-ottimizzata sul posto --- non come uno stato riciclato da
un punto diverso.

\messaggio{Due errori del progetto, misurati separatamente per la prima
volta: quello di preparazione (quanto bene l'ansatz raggiunge il
fondamentale) e quello di evoluzione (quanto bene Trotter approssima
$e^{-iHt}$). Non sono intercambiabili, e qui si vede quale dei due
conviene attaccare per primo.}

\section{I correlatori}

Correlatori dinamici $C_{ij}^{\alpha\beta}(t)=\media{\sigma_i^\alpha(t)\,
\sigma_j^\beta(0)}$, misurati con un circuito a tre qubit (due del dimero
più un'ancilla, schema Hadamard test). Punto di lavoro ``test 2''
($b/J=0.35$, $D/J=0.80$) --- lo stesso punto singolo già usato come
verifica aggiuntiva nel passo del VQE, non il punto della dinamica.

Perché proprio questo punto, e non uno dei due già incontrati: la
scansione in campo (Documenti 1--2) serve a vedere l'incrocio, quindi
copre un intervallo di $b/J$, non un punto singolo rappresentativo --- non
avrebbe senso per un correlatore, che va calcolato in un punto preciso.
Il punto dinamico ($b/J=-0.18$, $D/J=1$) è scelto per dare una dinamica
con più frequenze visibili, una condizione su come uno stato
\emph{fissato} risponde nel tempo --- non sulla natura dello stato
fondamentale in sé. ``Test 2'' è invece scelto apposta perché il
fondamentale \emph{stesso} sia non banale (una sovrapposizione genuina fra
settori, non uno stato semplice) e il DM sia ben acceso: è la condizione
giusta perché il confronto fra ansatz --- il cuore di questo passo quanto
di quello del VQE --- abbia qualcosa di non ovvio da mostrare. Le due
condizioni (dinamica ricca; fondamentale non banale) non coincidono per
forza nello stesso punto, ed è per questo che restano due punti distinti,
non uno.

\begin{center}
\small
\renewcommand{\arraystretch}{1.35}
\begin{tabular}{@{}p{2.6cm}p{3.2cm}p{7.6cm}@{}}
\toprule
\rowcolor{grigetto}
\textbf{Punto} & \textbf{Parametri} & \textbf{Perché sì o perché no, qui} \\
\midrule
Scansione in campo &
$D=0$, $D/J=0.2$, $b/J$ variabile &
No: copre un intervallo, non è un punto singolo --- un correlatore va
calcolato in un punto preciso \\
Punto dinamico &
$b/J=-0.18$, $D/J=1$ &
No: scelto per dinamica ricca (proprietà di uno stato \emph{fissato} nel
tempo), non per la natura del fondamentale \\
``Test 2'' \textit{(usato qui)} &
$b/J=0.35$, $D/J=0.80$ &
Sì: fondamentale stesso non banale, DM ben acceso --- la condizione
giusta per un confronto fra ansatz interessante \\
\bottomrule
\end{tabular}
\end{center}

Il grosso del passo usa il fondamentale esatto; solo nella sezione finale si
ripete l'operazione già vista nella dinamica --- e con la stessa
distinzione fra i due ansatz, non uno stato generico ``del VQE'':

\begin{itemize}[leftmargin=1.8em,itemsep=3pt]
\item un correlatore singolo, confrontato preparandolo con \textbf{entrambi}
gli ansatz (PMA base e PMA esteso, stessa struttura di circuito della
dinamica, ri-ottimizzata a questo punto), affiancati in due pannelli
--- stesso schema già visto;
\item uno scan sistematico sulle $36$ combinazioni, con \textbf{entrambi}
gli ansatz e \textbf{entrambi} i metodi di calcolo --- \emph{esatto}
($e^{-iHt}$ diretto, anche sullo stato imperfetto del VQE) e \emph{vero}
(circuito ad ancilla, $N=200$ passi di Trotter transpilati) --- tenuti
separati, non l'uno al posto dell'altro, e confrontati alla fine. Griglie
temporali diverse per i due metodi ($161$ punti l'esatto, $21$ il vero,
per tempo di calcolo) --- verificato che non cambi la conclusione:
infittendo a $81$ punti sulla combinazione peggiore, lo scarto si sposta
da $0.641$ a $0.646$, sotto l'$1\%$.
\end{itemize}

Il documento presenta le due sezioni di metodo una dopo l'altra, ciascuna
con le proprie figure e le proprie osservazioni, seguite da un confronto
esplicito --- non una sostituzione silenziosa di un metodo con l'altro.

Con il metodo \emph{esatto}, il PMA base mostra un fattore $\sim11\times$
fra combinazione più protetta e più fragile (pattern netto); il PMA
esteso dà scarti uniformemente minuscoli ($10^{-7}$--$10^{-8}$), senza
nulla di fisico da leggere --- quel calcolo confronta due evoluzioni
\emph{entrambe} esatte, un confronto che nessun circuito reale può
realizzare.

Con il metodo \emph{vero} (il circuito), il PMA base conferma lo stesso
pattern ($\sim12\times$, scarto fra $0.053$ e $0.64$) --- la struttura del
blocco CNOT--$R_y$--CNOT protegge la direzione $y$. Il PMA esteso, però,
\textbf{non} conferma il risultato del metodo esatto: lo scarto è
$0.011$--$0.032$, quattro-cinque ordini di grandezza più grande del
rumore numerico trovato prima --- il vero errore di Trotter, piccolo ma
reale, che il metodo esatto non può mostrare per costruzione.

\messaggio{Per il PMA base i due metodi concordano: l'errore di
preparazione domina, e qualunque modo si scelga di calcolare l'evoluzione
non cambia la conclusione. Per il PMA esteso, dove l'errore di
preparazione è quasi assente, il metodo di calcolo diventa la fonte
principale dell'errore residuo --- e un metodo che assume un'evoluzione
impossibile su hardware reale nasconde esattamente l'errore che il
circuito vero rivela. Scegliere un solo correlatore, o un solo metodo,
per riportare un risultato senza dichiararlo sarebbe fuorviante in
entrambi i casi.}

\clearpage
\section{Schema visivo d'insieme}

\begin{figure}[htbp]
\centering
\begin{tikzpicture}[
  box/.style={draw,thick,rounded corners=4pt,align=left,
              font=\footnotesize,inner sep=8pt,
              text width=5.6cm,minimum height=3.2cm},
  lab/.style={font=\scriptsize\itshape,text=gray!70!black,align=center},
  arrow/.style={-{Stealth[length=8pt]},thick},
  darrow/.style={-{Stealth[length=8pt]},thick,dashed}
]

\node[box,draw=blu,fill=blu!5] (sis) at (0,4.6) {
  \textbf{\textcolor{blu}{Il sistema}}\\[3pt]
  spettro esatto, nessun circuito\\
  $D{=}0$ vs $D/J{=}0.2$\\[3pt]
  \textit{scopre:} il DM rompe la\\
  conservazione di $M$
};

\node[box,draw=arancio,fill=arancio!7] (vqe) at (9.2,4.6) {
  \textbf{\textcolor{arancio!80!black}{Il VQE}}\\[3pt]
  sweep $b/J$, $D{=}0$/$0.2$\\
  $+$ punto ``test 2''\\[3pt]
  \textit{scopre:} ansatz-$M$ crolla\\
  con $D{\neq}0$; RBS$+2$ lo ripara
};

\node[box,draw=verdes,fill=verdes!6] (din) at (2.3,-1.0) {
  \textbf{\textcolor{verdes}{La dinamica}}\\[3pt]
  $b/J{=}-0.18$,\ $D/J{=}1$\\
  stato: $\ket{\uparrow\uparrow}$\\[3pt]
  \textit{scopre:} errore preparazione\\
  $10\times$ errore Trotter
};

\node[box,draw=violas,fill=violas!8] (cor) at (9.2,-1.0) {
  \textbf{\textcolor{violas}{I correlatori}}\\[3pt]
  punto ``test 2''\\
  stato: fondamentale esatto\\[3pt]
  \textit{scopre:} l'effetto dipende\\
  dalla combinazione di spin ($10\times$)
};

\draw[arrow,blu] (sis) -- (vqe) node[midway,above,lab] {motiva};

% un unico punto di biforcazione sotto il VQE, con l'etichetta condivisa
% messa lì una volta sola, non ripetuta su ciascun ramo (altrimenti le due
% etichette si sovrappongono dove i rami sono vicini)
\coordinate (fork) at ($(vqe.south)+(0,-1.1)$);
\draw[darrow,arancio!70!black] (vqe.south) -- (fork);
\node[lab,align=center,text width=4.6cm] at ($(fork)+(1.9,0.15)$)
  {stessa struttura di circuito,\\ri-ottimizzata a ogni punto};
\draw[darrow,arancio!70!black] (fork) -- (din.north east);
\draw[darrow,arancio!70!black] (fork) -- (cor.north);

\end{tikzpicture}
\caption{Frecce continue: collegamento diretto (un risultato motiva il
passo successivo). Frecce tratteggiate: riuso della \emph{struttura di
circuito} del VQE, ri-ottimizzata da zero al punto di lavoro proprio di
ciascun passo --- non un dato che passa direttamente da un passo
all'altro. La dinamica e i correlatori non sono collegati fra loro: sono
due applicazioni indipendenti della stessa idea, a punti diversi.}
\end{figure}

\end{document}
